\documentclass[11pt,a4paper]{article}

% ─── Packages ────────────────────────────────────────────────────────────────
\usepackage[top=2.2cm,bottom=2.2cm,left=2.2cm,right=2.2cm]{geometry}
\usepackage{amsmath, amssymb}
\usepackage{graphicx}
\usepackage{xcolor}
\usepackage{listings}
\usepackage{caption}
\usepackage{booktabs}
\usepackage{array}
\usepackage{enumitem}
\usepackage{fancyhdr}
\usepackage{float}
\usepackage{hyperref}
\usepackage{microtype}

% ─── Hyperref ────────────────────────────────────────────────────────────────
\hypersetup{colorlinks=true, linkcolor=blue!60!black, citecolor=green!50!black}

% ─── Page Style ──────────────────────────────────────────────────────────────
\pagestyle{fancy}
\fancyhf{}
\lhead{\small Tiny-ImageNet-10 Experiments}
\rhead{\small Problem 2.1 --- The Latent Walk (Interpolation)}
\cfoot{\thepage}
\renewcommand{\headrulewidth}{0.4pt}

% ─── Code Style ──────────────────────────────────────────────────────────────
\definecolor{codegreen}{RGB}{0,120,0}
\definecolor{codegray}{RGB}{110,110,110}
\definecolor{codeblue}{RGB}{0,0,200}
\definecolor{backcolor}{RGB}{248,248,248}

\lstdefinestyle{pystyle}{
    backgroundcolor=\color{backcolor},
    commentstyle=\color{codegreen},
    keywordstyle=\color{codeblue}\bfseries,
    stringstyle=\color{codegray},
    basicstyle=\ttfamily\footnotesize,
    breaklines=true,
    numbers=left,
    numberstyle=\tiny\color{codegray},
    numbersep=7pt,
    keepspaces=true,
    frame=single,
    language=Python,
    showstringspaces=false,
    tabsize=4
}
\lstset{style=pystyle}

% ─── Paragraph spacing ───────────────────────────────────────────────────────
\setlength{\parskip}{5pt}
\setlength{\parindent}{0pt}

% ─────────────────────────────────────────────────────────────────────────────
\begin{document}

% ─── Title Block (matching Problem 1.2 style) ────────────────────────────────
\begin{center}
    \rule{\linewidth}{1.5pt}\\[5pt]
    {\large\bfseries Problem 2.1 --- The Latent Walk (Interpolation)}\\[4pt]
    \rule{\linewidth}{0.8pt}
\end{center}
\thispagestyle{fancy}

% ══════════════════════════════════════════════════════════════════════════════
\subsection*{Background and Objective}
% ══════════════════════════════════════════════════════════════════════════════
\noindent\rule{\linewidth}{0.4pt}

A core claim of the Variational Autoencoder (VAE) is that it learns a \textbf{structured and
continuous latent space}. Unlike a deterministic autoencoder whose encoder can place codes
freely in $\mathbb{R}^{d}$, the VAE imposes a Gaussian prior through the KL-Divergence
regulariser. This forces nearby classes to occupy overlapping regions of $\mathcal{Z}$, and
crucially ensures the \emph{space between} encoded images is also decodable into semantically
meaningful outputs.

The \textbf{latent walk} is the standard empirical test of this property. Given two images from
different classes, we encode them, linearly interpolate their latent vectors in 10 steps of
$\alpha$, and decode each intermediate vector. The resulting transition strip directly reveals
whether the manifold is smooth (VAE-style) or whether intermediate points decode to incoherent
ghostly overlaps (standard AE/FCNN-style artefact).

% ══════════════════════════════════════════════════════════════════════════════
\subsection*{Loss Function Formulation}
% ══════════════════════════════════════════════════════════════════════════════
\noindent\rule{\linewidth}{0.4pt}

The CVAE minimises the negative ELBO, decomposed as:
\[
    \mathcal{L}_{\text{total}} = \mathcal{L}_{\text{recon}} + \beta \cdot \mathcal{L}_{\text{KLD}}
\]

\textbf{Reconstruction Loss --- Mean Squared Error (MSE)}
\[
    \mathcal{L}_{\text{recon}}
    = \frac{1}{N}\sum_{i=1}^{N}\sum_{d=1}^{D}
      \bigl(x_{id} - \hat{x}_{id}\bigr)^{2}
\]
$x_{id}$: true pixel, $\hat{x}_{id}$: decoder output,
$N$: batch size, $D = 3\times64\times64 = 12{,}288$.

\textbf{KL Divergence}
\[
    \mathcal{L}_{\text{KLD}}
    = -\frac{1}{2N}\sum_{i=1}^{N}\sum_{j=1}^{d}
      \bigl(1 + \log\sigma_{ij}^{2} - \mu_{ij}^{2} - \sigma_{ij}^{2}\bigr)
\]

\textbf{Interpolation Formula}

Given posterior means $\mathbf{z}_1 = \boldsymbol{\mu}_A$ and $\mathbf{z}_2 = \boldsymbol{\mu}_B$
for source images $\mathbf{x}_A$ and $\mathbf{x}_B$:
\[
    \mathbf{z}_{\text{new}} = \alpha\,\mathbf{z}_1 + (1-\alpha)\,\mathbf{z}_2,
    \qquad \alpha \in \{0.0,\;0.1,\;\ldots,\;1.0\}
\]
Posterior \emph{means} (not stochastic samples) are used so the walk is deterministic and reproducible across runs.

% ══════════════════════════════════════════════════════════════════════════════
\subsection*{Model Architecture}
% ══════════════════════════════════════════════════════════════════════════════
\noindent\rule{\linewidth}{0.4pt}

Convolutional encoder--decoder on $3\times64\times64$ RGB images.
Latent dim: 128. Optimiser: Adam ($10^{-3}$). Batch: 64. $\beta = 1.0$.

\begin{center}
\begin{tabular}{lll}
\toprule
\textbf{Component} & \textbf{Layer} & \textbf{Dimensions} \\
\midrule
Encoder   & Conv2d $\times$4, stride 2, ReLU      & $3\!\times\!64^2 \to 32\!\times\!32^2 \to 64\!\times\!16^2 \to 128\!\times\!8^2 \to 256\!\times\!4^2$ \\
          & Flatten                               & $4{,}096$ \\
          & FC (mean) / FC (log-var)              & $4{,}096 \to 128$ each \\
Reparam.  & $\mathbf{z}=\boldsymbol{\mu}+\boldsymbol{\varepsilon}\odot\boldsymbol{\sigma},\;\boldsymbol{\varepsilon}\sim\mathcal{N}(\mathbf{0},\mathbf{I})$ & dim 128 \\
Decoder   & FC + ReLU, Unflatten                  & $128 \to 4{,}096 \to 256\!\times\!4^2$ \\
          & ConvTranspose2d $\times$3, ReLU        & $256\!\times\!4^2 \to 128\!\times\!8^2 \to 64\!\times\!16^2 \to 32\!\times\!32^2$ \\
          & ConvTranspose2d + \textbf{Sigmoid}    & $32\!\times\!32^2 \to 3\!\times\!64^2$ \\
\bottomrule
\end{tabular}
\end{center}

The final \texttt{Sigmoid} constrains decoder outputs to $[0,1]$, matching normalised inputs and the MSE loss scale.

% ══════════════════════════════════════════════════════════════════════════════
\subsection*{Results --- Training ($\beta = 1.0$, 20 Epochs)}
% ══════════════════════════════════════════════════════════════════════════════
\noindent\rule{\linewidth}{0.4pt}

\begin{figure}[H]
    \centering
    \includegraphics[width=0.88\textwidth]{vae_loss_plot.png}
    \caption{Total, Reconstruction (MSE), and KL Divergence losses over 20 epochs ($\beta = 1.0$).}
    \label{fig:loss}
\end{figure}

\begin{center}
\begin{tabular}{ccccp{5.2cm}}
\toprule
\textbf{Epoch} & \textbf{Total} & \textbf{Recon} & \textbf{KLD} & \textbf{Notes}\\
\midrule
1  & 660 & 643 & 17 & Fast initial drop; encoder not yet structured \\
5  & 430 & 395 & 35 & KLD rising; encoder learning posteriors \\
10 & 360 & 320 & 40 & Slow convergence; KLD approaching plateau \\
20 & 325 & 283 & 43 & Final values; stable equilibrium \\
\bottomrule
\end{tabular}
\end{center}

\textbf{Key observations:}
\begin{itemize}[leftmargin=*, itemsep=1pt]
    \item Total and Recon losses drop sharply in epochs 1--3 as the network learns coarse image
          structure, then converge slowly.
    \item KLD \emph{rises} from $\approx 17$ to a plateau near $43$ --- the encoder gradually
          invests in informative latent codes as reconstruction pressure increases.
          This rising-then-stabilising pattern indicates a healthy encoder--decoder equilibrium,
          not posterior collapse.
    \item The gap between Total and Recon ($\approx 43$ at epoch 20) equals the stabilised KLD,
          confirming the loss decomposition is self-consistent.
\end{itemize}

% ══════════════════════════════════════════════════════════════════════════════
\subsection*{Results --- Latent Space Walk}
% ══════════════════════════════════════════════════════════════════════════════
\noindent\rule{\linewidth}{0.4pt}

\begin{figure}[H]
    \centering
    \includegraphics[width=\textwidth]{latent_walk_strip.png}
    \caption{Ten-step latent walk from Class A ($\alpha = 0.0$, dark textured landscape) to
             Class B ($\alpha = 1.0$, bright sky-dominant scene). Each frame is decoded from
             $\mathbf{z}_{\text{new}} = \alpha\boldsymbol{\mu}_A + (1-\alpha)\boldsymbol{\mu}_B$.}
    \label{fig:walk}
\end{figure}

\begin{center}
\begin{tabular}{p{2.4cm}p{5.4cm}p{5.4cm}}
\toprule
\textbf{$\alpha$ range} & \textbf{Visual appearance} & \textbf{Interpretation} \\
\midrule
$0.0$--$0.2$ & Dark olive/brown tones; textured foreground   & $\mathbf{x}_A$-dominant: low-light landscape \\
$0.3$--$0.4$ & Gradual lightening; horizon becomes visible   & Structural blending begins \\
$0.5$        & Equal brightness top/bottom; horizon centred  & Mid-latent: mixed structural identity \\
$0.6$--$0.7$ & Lighter; sky palette emerging above horizon   & $\mathbf{x}_B$ features dominate \\
$0.8$--$1.0$ & Grey-blue sky with dark horizon band          & $\mathbf{x}_B$-dominant: bright sky scene \\
\bottomrule
\end{tabular}
\end{center}

% ══════════════════════════════════════════════════════════════════════════════
\subsection*{Observation and Analysis}
% ══════════════════════════════════════════════════════════════════════════════
\noindent\rule{\linewidth}{0.4pt}

\textbf{Does the transition look like a ghostly overlap (FCNN-style) or a smooth structural
change (VAE-style)?}

\begin{center}
\begin{tabular}{p{3.8cm}p{5.0cm}p{5.0cm}}
\toprule
\textbf{Property} & \textbf{FCNN-style (undesirable)} & \textbf{VAE-style (observed)} \\
\midrule
Intermediate frames   & Both source images simultaneously visible as a double-exposure
                      & Gradual structural morphing; no double-exposure \\
Colour palette        & Inconsistent, abrupt hue jumps
                      & Monotonic shift: warm dark $\to$ cool bright \\
Structural continuity & Discontinuous feature boundaries
                      & Horizon migrates smoothly across all 10 frames \\
Semantic coherence    & Nonsensical compositions
                      & Physically plausible ``in-between'' scenes \\
\bottomrule
\end{tabular}
\end{center}

The observed transition is \textbf{VAE-style}: the horizon line migrates continuously across the
strip; luminance changes monotonically; no ghostly overlay of both source images appears at
any $\alpha$.

\textbf{What happens to sharpness at intermediate $\alpha$?}

\begin{center}
\begin{tabular}{p{3cm}p{4cm}p{6.5cm}}
\toprule
$\alpha$ & Sharpness & Reason \\
\midrule
$0.0$ or $1.0$ & Sharpest (endpoints) & Direct encoding of a training image \\
$0.2$--$0.8$   & Progressively blurrier toward $\alpha=0.5$ & Latent code lies between two training modes \\
$0.5$          & Blurriest (midpoint)  & Decoder has seen fewest examples near this code \\
\bottomrule
\end{tabular}
\end{center}

Blurriness at intermediate steps is an inherent VAE property caused by two compounding factors:

\textbf{MSE reconstruction loss.} Minimising squared error is equivalent to maximising a Gaussian
likelihood, which outputs the \emph{mean} of all plausible reconstructions --- inherently blurry.

\textbf{Posterior averaging.} Codes $\mathbf{z}_{\text{new}}(\alpha)$ occupy regions between two
training-distribution modes. The decoder sees fewer examples near these points, defaulting to a
conservative blurry average. This is not a failure of the walk; it is the fundamental
reconstruction--regularity trade-off of the VAE framework.

\textbf{Why does the KLD term make smooth walks possible?}

Without KLD, the encoder maps each input to an isolated point in $\mathbb{R}^{128}$; the decoder
is trained only at those points and produces noise everywhere else. Adding the KLD term forces all
posteriors toward $\mathcal{N}(\mathbf{0},\mathbf{I})$:
\[
    D_{\mathrm{KL}}\!\bigl(\mathcal{N}(\boldsymbol{\mu},\,\mathrm{diag}(\boldsymbol{\sigma}^2))
    \,\|\,\mathcal{N}(\mathbf{0},\mathbf{I})\bigr)
    = \frac{1}{2}\sum_{j=1}^{128}\bigl(\mu_j^2 + \sigma_j^2 - 1 - \log\sigma_j^2\bigr)
\]
This achieves two structural effects: (1)~\textbf{clustering} --- all class codes are pulled
toward the origin, reducing inter-class distances in $\mathcal{Z}$; (2)~\textbf{overlap} ---
non-zero $\sigma$ spreads each posterior over a neighbourhood, compelling the decoder to produce
valid images across a region, not just at a single code point. Together, these guarantee the
interpolation path passes through decoder-covered territory, making the walk coherent.

% ══════════════════════════════════════════════════════════════════════════════
\subsection*{Implementation}
% ══════════════════════════════════════════════════════════════════════════════
\noindent\rule{\linewidth}{0.4pt}

\begin{lstlisting}[caption={Latent walk: encode with mu, interpolate, decode.}]
def generate_latent_walk(encoder, decoder, image_1, image_2, device='cpu'):
    encoder.eval(); decoder.eval()
    image_1 = image_1.unsqueeze(0).to(device)
    image_2 = image_2.unsqueeze(0).to(device)

    with torch.no_grad():
        mu_1, _ = encoder(image_1)   # use mean (mu), not a stochastic sample
        mu_2, _ = encoder(image_2)
        z_1, z_2 = mu_1.squeeze(), mu_2.squeeze()

        alphas = torch.linspace(0, 1, steps=10).to(device)
        decoded = []
        for alpha in alphas:
            z_new = alpha * z_1 + (1 - alpha) * z_2   # linear interpolation
            decoded.append(decoder(z_new.unsqueeze(0)).squeeze().cpu())

    fig, axes = plt.subplots(1, 10, figsize=(20, 3))
    for i, ax in enumerate(axes):
        ax.imshow(np.clip(decoded[i].permute(1, 2, 0).numpy(), 0, 1))
        ax.axis('off')
        ax.set_title(f'a = {alphas[i]:.1f}')
    plt.suptitle("Latent Space Walk: Class A to Class B", fontsize=14)
    plt.tight_layout(); plt.show()
\end{lstlisting}

% ══════════════════════════════════════════════════════════════════════════════
\subsection*{Summary}
% ══════════════════════════════════════════════════════════════════════════════
\noindent\rule{\linewidth}{0.4pt}

\begin{enumerate}[leftmargin=*, itemsep=2pt]
    \item Latent walk uses $\mathbf{z}_{\text{new}} = \alpha\boldsymbol{\mu}_A +
          (1-\alpha)\boldsymbol{\mu}_B$ for 10 values of $\alpha$ from 0 to 1.
    \item At $\beta = 1.0$ over 20 epochs: total loss $660 \to 325$;
          KLD rises $17 \to 43$; Recon $643 \to 283$.
    \item The transition strip shows a \textbf{smooth structural change} (VAE-style):
          the horizon migrates continuously, luminance shifts monotonically, and
          no ghostly double-exposure artefact appears.
    \item Higher $\alpha$ blurriness near the midpoint is an expected VAE property
          caused by MSE loss and posterior averaging --- not a failure of the interpolation.
    \item The KLD term is the enabling mechanism: it clusters and overlaps posterior
          distributions so the decoder covers inter-class regions of $\mathcal{Z}$.
\end{enumerate}

% ══════════════════════════════════════════════════════════════════════════════
\subsection*{References}
% ══════════════════════════════════════════════════════════════════════════════
\noindent\rule{\linewidth}{0.4pt}

\begin{enumerate}[label={[\arabic*]}, leftmargin=*, itemsep=1pt]
    \item D.\,P. Kingma and M. Welling, ``Auto-Encoding Variational Bayes,'' \textit{ICLR}, 2014.
    \item I. Higgins et al., ``$\beta$-VAE: Learning Basic Visual Concepts with a Constrained
          Variational Framework,'' \textit{ICLR}, 2017.
    \item A. Radford, L. Metz, and S. Chintala, ``Unsupervised Representation Learning with
          Deep Convolutional Generative Adversarial Networks,'' \textit{ICLR}, 2016.
    \item Le et al., ``Tiny ImageNet Visual Recognition Challenge,'' \textit{CS 231N}, 2015.
\end{enumerate}

\end{document}
